\chapter{Conditions}\label{chap:conditions}

Condition evaluation maps condition blocks to three-valued results.
The key invariant: the no-context evaluation never hides a grant
visible under any concrete context.

\section{Three-Valued Logic}

\begin{definition}[Tri]\label{def:Tri}
\lean{Tri}
\leanok
Three-valued type: \texttt{.t} (true), \texttt{.f} (false),
\texttt{.u} (unknown).
\end{definition}

\begin{definition}[Tri.and]\label{def:Tri.and}
\lean{Tri.and}
\leanok
Kleene conjunction.
\end{definition}

\begin{definition}[Tri.or]\label{def:Tri.or}
\lean{Tri.or}
\leanok
Kleene disjunction.
\end{definition}

\begin{definition}[Tri.not]\label{def:Tri.not}
\lean{Tri.not}
\leanok
Kleene negation.
\end{definition}

\begin{lemma}[Tri.and\_tu]\label{lem:Tri.and_tu}
\lean{Tri.and_tu}
\leanok
If both operands are T or U, then \texttt{Tri.and} produces T or U.
\end{lemma}

\begin{lemma}[Tri.and\_eq\_t]\label{lem:Tri.and_eq_t}
\lean{Tri.and_eq_t}
\leanok
\texttt{Tri.and a b = .t} iff \texttt{a = .t} and \texttt{b = .t}.
\end{lemma}

\section{Condition Evaluation}

\begin{definition}[evalCond]\label{def:evalCond}
\lean{evalCond}
\leanok
Evaluates a list of condition blocks against a context, producing
a \texttt{Tri} result and a list of warnings.  Blocks are conjoined
(\texttt{Tri.and}).
\end{definition}

\begin{theorem}[evalCond\_noContext\_tu]\label{thm:evalCond_noContext_tu}
\uses{lem:Tri.and_tu}
\lean{evalCond_noContext_tu}
\leanok
Under \texttt{noContext}, every condition block evaluates to T or U,
never F.  This is the foundation of the grant-preserving semantics:
no-context never suppresses an Allow.
\end{theorem}

\begin{theorem}[evalCond\_noContext\_T\_imp]\label{thm:evalCond_noContext_T_imp}
\uses{lem:Tri.and_eq_t}
\lean{evalCond_noContext_T_imp}
\leanok
If a condition block evaluates to T under \texttt{noContext}, it
evaluates to T under every context.  This ensures Deny applicability
is preserved under no-context.
\end{theorem}

\section{No-Context Conservatism}

\begin{corollary}[allows\_nocontext\_conservative]\label{cor:allows_nocontext_conservative}
\uses{thm:nocontext_conservative}
\lean{allows_nocontext_conservative}
\leanok
If an IAM policy allows a request under any context, it allows the
same request under \texttt{noContext}.  The no-context report never
hides a grant that some complete context would allow.  Direct
application of the generic \texttt{nocontext\_conservative} via
\texttt{iamSem}.
\end{corollary}
